\documentclass[11pt,a4paper]{article}
\RequirePackage{fontspec}
\usepackage[ngerman]{babel}
\defaultfontfeatures{Ligatures=TeX}
\usepackage{amsmath,amssymb}
\usepackage{geometry}
\geometry{left=25mm, right=25mm, top=25mm, bottom=25mm}
\usepackage{booktabs}
\usepackage{array}
\usepackage{tabularx}
\usepackage{xcolor}
\usepackage{float}
\usepackage{enumitem}
\usepackage{fancyhdr}
\usepackage{titlesec}
\usepackage{hyperref}
\usepackage{siunitx}
\usepackage{tcolorbox}

\definecolor{darkblue}{HTML}{16213e}
\definecolor{accentred}{HTML}{e94560}
\definecolor{keygreen}{HTML}{2e7d32}
\definecolor{warnred}{HTML}{c62828}

\newtcolorbox{keybox}{colback=keygreen!8, colframe=keygreen!60!black, 
  left=3mm, right=3mm, top=2mm, bottom=2mm, fonttitle=\bfseries}
\newtcolorbox{warnbox}{colback=warnred!8, colframe=warnred!60!black,
  left=3mm, right=3mm, top=2mm, bottom=2mm, fonttitle=\bfseries}

\pagestyle{fancy}
\fancyhf{}
\fancyhead[L]{\small Dok.\,0005 -- Frequenzverschiebung als Indikator der Ansprache}
\fancyhead[R]{\small\thepage}
\fancyfoot[C]{\small Querverweise: Dok.\,0002 (v8) $\cdot$ Dok.\,0003 (Impedanzvergleich) $\cdot$ Dok.\,0004 (Zwei-Filter)}

\titleformat{\section}{\Large\bfseries\color{darkblue}}{Kapitel \thesection:}{0.5em}{}
\titleformat{\subsection}{\large\bfseries\color{darkblue!80}}{}{0em}{}

\begin{document}

\begin{center}
{\LARGE\bfseries\color{darkblue} Dok.\,0005: Frequenzverschiebung als Indikator der Ansprache}\\[6pt]
{\large Beweis des Zusammenhangs: $\Delta f$ und $\Delta\tau$ als Real- und Imaginärteil derselben Kopplung}\\[4pt]
\textcolor{accentred}{\rule{0.8\textwidth}{2pt}}
\end{center}
{\small\itshape Beweist, dass Frequenzverschiebung und Ansprache-Verschlechterung physikalisch untrennbar verknüpft sind. Beide sind Projektionen derselben komplexen Kopplungsimpedanz $Z_H(f)$. Zahlenwerte: 50-Hz-Basszunge aus Dok.\,0002.}

\vspace{2mm}
\begin{center}\small
\begin{tabular}{ll}
\toprule
\textbf{Dokumentenverweise} \\
\midrule
Dok.\,0001 & Analyse zur Instrumentenakustik und Wahrnehmungspsychologie \\
Dok.\,0002 & Strömungsanalyse Bass-Stimmzunge 50\,Hz -- v8 \\
Dok.\,0003 & Impedanzvergleich Durchschlagzunge vs.\ Labialpfeife \\
Dok.\,0004 & Frequenzvariation -- Zwei-Filter-Modell \\
\textbf{Dok.\,0005} & \textbf{Frequenzverschiebung als Indikator der Ansprache (dieses Dokument)} \\
Dok.\,0006 & Zeichenerklärung \\
Dok.\,0007 & Diskant-Stimmstock -- Kammerfrequenzen \\
Dok.\,0008 & Klangveränderung durch Kammergeometrie \\
Dok.\,0500 & Leitfaden zur ästhetischen Forensik bei Akkordeon-Gehäusen \\
\midrule
\href{berechnungen_verifikation.py}{berechnungen\_verifikation.py} & Verifikationsskript aller Berechnungen \\
\bottomrule
\end{tabular}
\end{center}
\vspace{4mm}

% ============================================================
\section{Die zentrale These}
% ============================================================

Dok.\,0002 (Kap.\,10) zeigt: Eine resonante Kammer verschlechtert die Ansprache über Serienkopplung. Dok.\,0004 zeigt: Die Kammer verschiebt die Zungenfrequenz um 1--5~Cent. Beide Effekte werden durch dieselbe physikalische Größe verursacht -- die \textbf{komplexe Kopplungsimpedanz} $Z_H(f)$.

\begin{keybox}
\textbf{Frequenzverschiebung und Ansprache-Verschlechterung sind der Imaginär- und Realteil derselben komplexen Kopplungsfunktion.} Sie sind durch die Kramers-Kronig-Relation kausal verknüpft. An der Frequenzverschiebung lässt sich ablesen, ob die Ansprache besser oder schlechter wird -- und wie weit man von der kritischen Resonanz entfernt ist.
\end{keybox}

% ============================================================
\section{Die komplexe Kopplungsimpedanz}
% ============================================================

\subsection{Zerlegung in Real- und Imaginärteil}

Die Helmholtz-Impedanz der Kammer bei Frequenz $f$:
\begin{equation}
Z_H(f) = R_H(f) + j\,X_H(f)
\end{equation}

\textbf{Realteil} (Resistanz $=$ Energiedissipation $\to$ Ansprache-Verschlechterung):
\begin{equation}
R_H(f) = Z_0 \cdot \frac{f/f_H}{Q_H \cdot D(f)}
\end{equation}

\textbf{Imaginärteil} (Reaktanz $=$ Energiespeicherung $\to$ Frequenzverschiebung):
\begin{equation}
X_H(f) = Z_0 \cdot \frac{1 - (f/f_H)^2}{D(f)}
\end{equation}

wobei $D(f) = \left(1 - (f/f_H)^2\right)^2 + \left(f/(f_H \cdot Q_H)\right)^2$ und $Z_0 = 1/(\omega_H C_a)$.

\subsection{Physikalische Zuordnung}

\textbf{Realteil $R_H(f) \to$ Energieverlust $\to$ Ansprache-Verschlechterung:}
\begin{equation}
\Delta\zeta_n = \kappa_\text{eff}^2 \cdot \frac{R_H(nf_1)}{Z_0}, \qquad \Delta\tau_n = \frac{\Delta\zeta_n}{\pi f_1}
\end{equation}

\textbf{Imaginärteil $X_H(f) \to$ Frequenzverschiebung:}
\begin{equation}
\Delta f_n = -\kappa_\text{eff}^2 \cdot f_n \cdot \frac{X_H(f_n)}{2 Z_0}
\end{equation}

\begin{keybox}
$\Delta\tau$ (Ansprache) $\propto R_H(f)$ und $\Delta f$ (Frequenz) $\propto X_H(f)$. Beide stammen aus \textbf{derselben} komplexen Impedanz $Z_H = R_H + jX_H$. Sie verhalten sich zueinander wie der Sinus und Kosinus eines Winkels -- kennt man den einen, kennt man den anderen.
\end{keybox}

% ============================================================
\section{Der Vorzeichenwechsel als Warnsignal}
% ============================================================

Die entscheidende Erkenntnis liegt im \textbf{Vorzeichenwechsel} der Frequenzverschiebung:

\begin{itemize}[leftmargin=5mm, itemsep=4pt]
\item $\Delta f > 0$ (Ton wird \textbf{höher}): Nächster OT liegt \textbf{unter} $f_H$. Kammer wirkt als Feder (Compliance). Ansprache ist gut, verschlechtert sich wenn $\Delta f$ wächst.
\item $\Delta f \to 0$ \textbf{mit Vorzeichenwechsel}: OT durchfährt $f_H$. \textbf{Worst Case.} Maximale Energieextraktion bei minimalem Frequenzsignal.
\item $\Delta f < 0$ (Ton wird \textbf{tiefer}): Nächster OT liegt \textbf{über} $f_H$. Kammer wirkt als Masse (Inertanz). Ansprache war schlecht, wird wieder besser.
\item $\Delta f \approx 0$ \textbf{ohne Vorzeichenwechsel}: Kein OT nahe $f_H$. \textbf{Best Case.}
\end{itemize}

\begin{warnbox}
\textbf{Das Paradoxon:} Exakt bei optimaler Resonanz ($f_H = n \cdot f_\text{Zunge}$) ist die Dämpfung maximal ($R_H = Q_H \cdot Z_0$), aber die Frequenzverschiebung \textbf{dieses Obertons} ist Null ($X_H = 0$). Am Resonanzpunkt wird alle Kopplungsenergie in Dissipation umgewandelt (Realteil), nichts bleibt für Reaktanz (Imaginärteil). Der Oberton wird nicht verschoben, sondern \textbf{gedämpft} -- das ist schlimmer.
\end{warnbox}

\subsection{Nyquist-Diagramm der Kopplung}

Im Nyquist-Diagramm ($R_H$ horizontal, $X_H$ vertikal) beschreibt die Impedanz einen \textbf{Kreis}. Jeder Oberton liegt auf einem Punkt dieses Kreises:

\begin{table}[H]
\centering\small
\begin{tabular}{rrrl}
\toprule
OT $n$ & $R_H/Z_0$ & $X_H/Z_0$ & Position auf dem Kreis \\
\midrule
1 (50\,Hz) & 0{,}001 & $+$1{,}000 & $\uparrow$ Oberer Rand (reine Compliance) \\
3 (150\,Hz) & 0{,}011 & $+$0{,}917 & $\uparrow$ Obere Hälfte \\
5 (250\,Hz) & 0{,}090 & $+$0{,}537 & $\nearrow$ Annäherung an Gipfel \\
$f_H$ (274) & 0{,}143 & $\approx 0$ & $\to$ \textbf{Gipfel} (maximale Dämpfung) \\
6 (300\,Hz) & 0{,}072 & $-$0{,}374 & $\searrow$ Unterhalb des Gipfels \\
8 (400\,Hz) & 0{,}017 & $-$0{,}655 & $\downarrow$ Untere Hälfte \\
\bottomrule
\end{tabular}
\caption{Nyquist-Positionen der Obertöne. Am Gipfel ($f_H$): max.\ Dämpfung, $X_H$ wechselt Vorzeichen.}
\end{table}

Der 5.~Oberton (250\,Hz) liegt in der oberen Kreishälfte ($X > 0$, $\Delta f > 0$), der 6.~Oberton (300\,Hz) in der unteren ($X < 0$, $\Delta f < 0$). Dazwischen, bei $f_H = 274$\,Hz, liegt der Kreisgipfel. Dass kein Oberton exakt auf den Gipfel trifft, erklärt die relativ gute Ansprache der 50-Hz-Zunge in dieser Kammer.

% ============================================================
\section{Der mathematische Beweis: Kramers-Kronig}
% ============================================================

Die Verknüpfung ist \textbf{physikalisch zwingend}. Jede kausale, lineare Antwortfunktion erfüllt die Kramers-Kronig-Relationen:
\begin{align}
X_H(\omega) &= -\frac{1}{\pi}\, \mathcal{P} \int_{-\infty}^{\infty} \frac{R_H(\omega')}{\omega' - \omega}\, d\omega' \\
R_H(\omega) &= +\frac{1}{\pi}\, \mathcal{P} \int_{-\infty}^{\infty} \frac{X_H(\omega')}{\omega' - \omega}\, d\omega'
\end{align}
($\mathcal{P}$ = Cauchyscher Hauptwert). \textbf{Kennt man den Frequenzverlauf der Frequenzverschiebung (alle $X_H(f)$), kann man daraus die Ansprache-Verschlechterung berechnen (alle $R_H(f)$) -- und umgekehrt.}

Für den Helmholtz-Resonator:
\begin{equation}
Z_H(f) = \frac{Z_0}{1 - (f/f_H)^2 + j\,f/(f_H \cdot Q_H)}
\end{equation}
Der Realteil (Lorentz-Linie, Glockenform) und der Imaginärteil (Dispersionskurve, S-Form mit Nulldurchgang bei $f_H$) sind durch die Hilbert-Transformation verknüpft -- genau wie Absorption und Brechungsindex in der Optik.

\begin{keybox}
Die Kramers-Kronig-Relation ist eine Folge der \textbf{Kausalität}. Die Kammer kann nicht ``wählen'', nur Frequenz zu verschieben ohne Energie zu dissipieren (oder umgekehrt). Beide Effekte sind zwei Seiten derselben Medaille. Ein System, das die Frequenz stark zieht, \textbf{muss} auch Energie dissipieren -- und umgekehrt (außer exakt bei Resonanz, wo $X = 0$).
\end{keybox}

% ============================================================
\section{Die praktische Messvorschrift}
% ============================================================

\textbf{Schritt 1:} Zungenfrequenz \textbf{frei} messen (Stimmplatte auf Prüfblock, keine Kammer). Ergebnis: $f_\text{frei}$.

\textbf{Schritt 2:} Zungenfrequenz \textbf{in der Kammer} messen (gleicher Balgdruck). Ergebnis: $f_\text{Kammer}$.

\textbf{Schritt 3:} $\Delta f = f_\text{Kammer} - f_\text{frei}$.

\textbf{Schritt 4:} Interpretation:

\begin{table}[H]
\centering\small
\begin{tabular}{l l l}
\toprule
\textbf{Messergebnis} & \textbf{Physik} & \textbf{Ansprache} \\
\midrule
$\Delta f \approx 0$, Ansprache gut & Kein OT nahe $f_H$ & \textbf{Optimal} \\
$\Delta f > 0$ (Ton höher) & Nächster OT unter $f_H$ & Gut, wird schlechter \\
$\Delta f$ maximal positiv & OT nähert sich $f_H$ & \textbf{Warnung} \\
$\Delta f \to 0$ (Vorzeichenwechsel) & OT durchfährt $f_H$ & \textbf{Worst Case} \\
$\Delta f < 0$ (Ton tiefer) & Nächster OT über $f_H$ & Besserung \\
$\Delta f \approx 0$, Ansprache gut & OT weit von $f_H$ & \textbf{Optimal} \\
\bottomrule
\end{tabular}
\caption{Interpretationsschlüssel: $\Delta f \to$ Ansprache-Prognose}
\end{table}

\textbf{Elegante Variante:} Verstellbare Kammer (verschiebbarer Boden). Volumen $V$ und damit $f_H$ stetig variieren. Wenn $\Delta f$ das Vorzeichen wechselt, hat man die Resonanz durchfahren -- dort ist die Ansprache am schlechtesten. Die Synchronizität von Vorzeichenwechsel und Ansprache-Minimum ist der \textbf{experimentelle Beweis} der Kramers-Kronig-Verknüpfung.

% ============================================================
\section{Gesamtbilanz: Summe über alle Obertöne}
% ============================================================

Der Gesamteffekt auf den Grundton ist die Summe über alle Obertöne:
\begin{equation}
\Delta f_\text{gesamt} = \sum_n \frac{a_n}{n} \cdot \Delta f_n, \qquad
\Delta\tau_\text{gesamt} = \sum_n a_n^2 \cdot \Delta\tau_n
\end{equation}

Das Vorzeichen der Netto-Frequenzverschiebung zeigt an, von welcher Seite der dominante Oberton an $f_H$ herankommt:

\begin{keybox}
\textbf{Faustregel:}\\
$\Delta f_\text{netto} > 0$ $\to$ Dominanter OT liegt \textbf{unter} $f_H$ $\to$ $f_H$ senken verschlechtert Ansprache.\\
$\Delta f_\text{netto} < 0$ $\to$ Dominanter OT liegt \textbf{über} $f_H$ $\to$ $f_H$ heben verschlechtert Ansprache.\\
$\Delta f_\text{netto} \approx 0$ $\to$ Entweder entkoppelt (gut) oder auf Resonanz (schlecht) $\to$ Ansprache-Test entscheidet.
\end{keybox}

% ============================================================
\section{Das Problem der fehlenden Referenzfrequenz}
% ============================================================

Die Messvorschrift ($f_\text{frei}$ vs.\ $f_\text{Kammer}$) hat eine fundamentale \textbf{praktische Einschränkung}, die in den Berechnungen nicht auftaucht, aber für den Instrumentenbauer entscheidend ist:

\begin{warnbox}
\textbf{Es gibt keine absolute Referenzfrequenz.} Die Frequenz $f_\text{frei}$ (Zunge auf dem Prüfblock, ohne Kammer) ist \textbf{selbst} eine setup-abhängige Größe. Der Prüfblock hat ein eigenes Volumen, eine eigene akustische Impedanz, eine eigene Kopplung. $f_\text{frei}$ ist nicht die ``wahre'' Eigenfrequenz der Zunge -- es ist die Frequenz der Zunge \textbf{in einem anderen akustischen Umfeld}.
\end{warnbox}

Die physikalische Eigenfrequenz im Vakuum ($f_\text{vak}$) ist nicht praktisch messbar -- die Zunge braucht Luft zum Schwingen (Bernoulli-Mechanismus). Selbst eine Zunge im offenen Raum hat eine akustische Umgebung: die freie Schallabstrahlung als schwache Belastung.

\subsection{Die drei ``Referenzfrequenzen'' des Akkordeonbauers}

\textbf{$f_\text{Prüfblock}$:} Zunge auf einem standardisierten Holzblock, angeblasen mit definiertem Druck. Am nächsten an $f_\text{vak}$, aber nicht identisch (der Block hat eine Compliance). Geschätzter Unterschied zu $f_\text{vak}$: 0{,}1--0{,}5~Cent.

\textbf{$f_\text{Kammer,offen}$:} Zunge in der Kammer montiert, Klappe offen. Die offene Klappe verändert die Kammerimpedanz gegenüber dem Spielbetrieb (niedrigere $Z$ am Klappenende $\to$ $f_H$ verschiebt sich).

\textbf{$f_\text{Kammer,Spiel}$:} Zunge in der Kammer, Klappe geschlossen, vom Balg angeblasen. Das ist der ``echte'' Betriebszustand -- aber hier ist die Frequenz bereits durch die volle Kammerkopplung verändert.

\subsection{Das Vorstimm-Setup als lokale Referenz}

Da keine der drei Referenzfrequenzen ``absolut'' ist, arbeitet der erfahrene Instrumentenbauer mit einer \textbf{lokalen Referenz}: Er definiert sein persönliches Vorstimm-Setup (bestimmter Prüfblock, bestimmter Blasdruck, bestimmte Raumtemperatur) als \textbf{Nullpunkt}. Alle Frequenzmessungen werden relativ zu diesem Setup interpretiert. Das Setup ist die ``Eichung'' des Stimmvorgangs.

Die $\Delta f$-Methode funktioniert trotzdem, weil sie eine \textbf{Differenz} misst: $\Delta f = f_\text{Kammer} - f_\text{Setup}$. Die absolute Frequenz spielt keine Rolle -- nur die Differenz und vor allem der \textbf{Vorzeichenwechsel} beim Variieren der Kammergeometrie. Der Vorzeichenwechsel ist setup-unabhängig, weil er eine Eigenschaft der Kammer-Zungen-Kopplung ist, nicht der absoluten Frequenz.

\subsection{Was die Erfahrung liefert, was die Berechnung nicht kann}

\textbf{1.\ Die Offset-Korrektur:} Wie viel Cent liegt $f_\text{Setup}$ über oder unter $f_\text{vak}$? Das weiß nur, wer jahrelang mit demselben Prüfblock gestimmt hat und die resultierenden Instrumente spielen gehört hat.

\textbf{2.\ Die Temperaturkompensation:} $f_\text{Zunge}$ ändert sich mit der Temperatur (E-Modul: $\sim$0{,}5\,Cent/°C). $c$ ändert sich ebenfalls ($+0{,}17\,\%$/°C $\to$ verändert $f_H$ und $\Delta f$). Der erfahrene Stimmer kompensiert beides intuitiv.

\textbf{3.\ Die Druckabhängigkeit:} Der Balgdruck verändert $h_\text{eff}$ (Dok.\,0002 Kap.\,9) und damit $S_\text{Spalt}$, $\kappa$ und $\Delta f$. Das Vorstimm-Setup definiert einen bestimmten Arbeitspunkt auf der Druckkurve.

\textbf{4.\ Die Materialstreuung:} 0{,}1\,mm Unterschied in der Aufbiegung $\to$ 7\,\% in $S_\text{eff}$ (Dok.\,0002 Kap.\,13), und damit unterschiedliche $\kappa$ und $\Delta f$. Die Berechnung gibt den Mittelwert -- die Erfahrung sagt, wie weit der Einzelfall abweichen kann.

\begin{keybox}
\textbf{Die Rolle der Berechnung im Stimmvorgang:}

Die Formeln aus Dok.\,0002--0005 ersetzen nicht die Erfahrung. Sie leisten etwas anderes: Sie erklären, \textbf{warum} die Erfahrungsregeln funktionieren. Sie sagen vorher, \textbf{in welche Richtung} eine Änderung wirkt. Sie identifizieren die \textbf{empfindlichen Parameter}. Sie warnen, \textbf{wann} man aufpassen muss (Vorzeichenwechsel $\to$ Resonanz).

Was die Berechnung \textbf{nicht} kann: den exakten Cent-Wert vorhersagen, die Materialstreuung kompensieren, das Setup des Instrumentenbauers kennen. Dafür braucht es den praktischen Test -- und die Erfahrung, die aus tausend Tests ein Handwerk macht.
\end{keybox}

% ============================================================
\section{Klarstellung: Ansprache vs.\ Lautstärke -- kein Trade-off bei der Kammerabstimmung}
% ============================================================

Ein naheliegendes Missverständnis muss ausgeräumt werden: Wenn optimale Ansprache bedeutet, dass die Kammer der Zunge möglichst \textbf{wenig} Energie entzieht -- reduziert dann gute Ansprache die maximal erzielbare Lautstärke?

\textbf{Nein.} Die Antwort ergibt sich direkt aus der Physik der Serienkopplung (Dok.\,0002 Kap.\,10). Zwei Effekte müssen sauber getrennt werden, weil sie physikalisch verschieden sind:

\subsection{Kammerabstimmung: Win-Win}

Wenn ein Oberton der Zunge exakt auf $f_H$ liegt, ist $R_H$ maximal -- die Kammer entzieht der Zunge Energie über die Serienkopplung. Dieser Energieverlust verschlechtert \textbf{beides gleichzeitig}:

\textbf{Ansprache:} Die entzogene Energie fehlt beim Amplitudenaufbau. Die Einschwingzeit $\tau$ verlängert sich um $\Delta\tau \propto R_H$ (Kap.\,3).

\textbf{Lautstärke:} Im eingeschwungenen Zustand fließt ein Teil der Schwingungsenergie \textbf{pro Zyklus} in die Kammer-Dissipation ab, statt in Schallabstrahlung umgewandelt zu werden. Die Gleichgewichtsamplitude sinkt: $A_\text{max} \propto 1/\sqrt{R_H}$.

Beide Effekte stammen aus derselben Quelle -- $R_H(f)$, dem Realteil der Kammerimpedanz. Wenn $R_H$ sinkt (kein Oberton auf $f_H$), verbessern sich \textbf{beide}: schnellere Ansprache \textbf{und} größere Amplitude. Bei der Kammerabstimmung gibt es keinen Trade-off. Optimale Abstimmung ist Win-Win.

\subsection{Zungen-Güte $Q_1$: Der echte Trade-off}

Der Trade-off zwischen Ansprache und Lautstärke existiert -- aber er liegt \textbf{nicht} in der Kammerabstimmung, sondern in der Zunge selbst:

\textbf{Hohes $Q_1$} (geringe Dämpfung): Die Zunge speichert viel Energie pro Zyklus $\to$ \textbf{lauter}. Aber der Amplitudenaufbau dauert viele Zyklen $\to$ \textbf{langsame Ansprache} ($\tau = Q_1/(\pi f)$).

\textbf{Niedriges $Q_1$} (starke Dämpfung): Wenige Zyklen zum Einschwingen $\to$ \textbf{schnelle Ansprache}. Aber die Gleichgewichtsamplitude ist geringer $\to$ \textbf{leiser}.

Dieser Trade-off wird durch Material, Einspannung und Aufbiegung bestimmt (Dok.\,0002 Kap.\,9) -- \textbf{nicht} durch die Kammergeometrie.

\begin{keybox}
\textbf{Zusammenfassung der Unterscheidung:}

\textbf{Kammerabstimmung} ($f_H$ relativ zu Obertönen, Dok.\,0005): Beeinflusst die \textbf{externe Dämpfung} durch die Kammer. Optimierung = $R_H$ minimieren = Win-Win (bessere Ansprache \textit{und} mehr Lautstärke).

\textbf{Zungengüte} $Q_1$ (Material, Einspannung, Dok.\,0002): Beeinflusst die \textbf{interne Dämpfung} der Zunge. Hier gibt es den echten Trade-off: mehr $Q_1$ = lauter aber langsamer, weniger $Q_1$ = schneller aber leiser.

Der Instrumentenbauer optimiert beides \textbf{unabhängig}: Die Kammer so abstimmen, dass sie möglichst wenig stört ($f_H$ zwischen Obertönen). Die Zunge so wählen, dass $Q_1$ zum musikalischen Kontext passt (Bass: eher niedrig für schnelle Ansprache; Diskant: eher hoch für Lautstärke und Sustain).
\end{keybox}

% ============================================================
\section{Zusammenfassung}
% ============================================================

\begin{keybox}
\textbf{Bewiesener Zusammenhang:} $\Delta f$ und $\Delta\tau$ sind Imaginär- und Realteil derselben komplexen Kopplungsimpedanz $Z_H(f)$, verknüpft durch Kramers-Kronig (Kausalität).

\textbf{Optimale Resonanz ergibt schlechteste Ansprache:} Bei $f_H = n \cdot f_\text{Zunge}$ ist $R_H$ maximal, $X_H = 0$. Maximum der Energieextraktion bei Minimum der Frequenzverschiebung.

\textbf{Die Frequenzverschiebung ist ein messbarer Indikator:}
Großes $|\Delta f|$ $\to$ starke Kopplung, aber nicht kritisch.
$\Delta f \to 0$ mit Vorzeichenwechsel $\to$ \textbf{schlechteste Ansprache}.
$\Delta f \to 0$ ohne Vorzeichenwechsel $\to$ \textbf{beste Ansprache}.

\textbf{Messvorschrift:} $f_\text{frei}$ vs.\ $f_\text{Kammer}$. Der Vorzeichenwechsel beim Variieren der Kammergeometrie markiert den Resonanzdurchgang.

\textbf{Fehlende absolute Referenz:} $f_\text{frei}$ ist setup-abhängig. Das Vorstimm-Setup (Prüfblock, Blasdruck, Temperatur) definiert den lokalen Nullpunkt. Der Vorzeichenwechsel von $\Delta f$ ist setup-unabhängig -- er ist die physikalische Observable.

\textbf{Kein Trade-off bei der Kammerabstimmung:} Optimale Abstimmung ($f_H$ zwischen Obertönen) verbessert Ansprache \textit{und} Lautstärke gleichzeitig, weil beide unter derselben Quelle ($R_H$) leiden. Der Trade-off zwischen Ansprache und Lautstärke liegt in der Zungengüte $Q_1$ (Material, Einspannung) -- nicht in der Kammer.

\textbf{Die Berechnung sagt, wo man suchen soll. Das Setup definiert den Nullpunkt. Die Erfahrung sagt, was man findet.}
\end{keybox}

\vspace{3mm}
\begin{center}
\textcolor{accentred}{\rule{0.6\textwidth}{1pt}}\\[4pt]
{\small\itshape Was man nicht messen kann, kann man nicht optimieren. Die Frequenzverschiebung ist die Messgröße -- das Setup ist die Eichung.}
\end{center}

\end{document}
